\section{결론}

\subsection{연구 요약}

본 논문은 HNDL 공격에 직접 노출되는 TLS 통신 구간을 대상으로, Classical ECC에서 Hybrid PQC를 거쳐 Advanced Hybrid PQC로의 단계적 전환을 자동화하는 11-Step DevSecOps 파이프라인을 설계·구현하였다.
제안 시스템은 TLS 설정 변경·보안 스캔·동적 검증·CBOM 생성·성능 검증·자동 롤백을 단일 파이프라인으로 통합하며, Policy as Code 기반으로 각 전환 단계의 정책 준수 여부를 자동으로 집행한다.

평가를 통해 세 가지 핵심 결과를 확인하였다.
첫째, PQC 도구 체인 통합 과정에서 발생하는 인프라 함정 6종에 대해 표준 SCA/SAST/SBOM 도구는 0/6을 사전 탐지한 반면, 자체 빌드는 5/6을 예방하고 매트릭스 CI는 6/6을 감지하였다.
둘째, CycloneDX~1.7 기반 CBOM Diff를 통해 Stage 전환 전후 암호 자산 변화를 BASELINE·IMPROVED·UNCHANGED·REGRESSED로 자동 판정하여 전환 진척도를 객관적으로 추적하는 구조를 설계·구현하였다.
셋째, LLM 기반 코드 마이그레이션 80 trial 정량 평가에서 정적 AST 검증 단독으로는 시그니처 환각 24건(44\%)을 차단하지 못하며, 의미적 실행 검증을 추가할 경우 본 평가 환경에서는 실패 코드가 배포 단계로 전달되지 않음을 확인하였다.
이는 LLM 마이그레이션의 성공을 의미하기보다는, 현재 모델 출력에 대해 런타임 검증이 필수적인 안전장치임을 보여준다.

\subsection{한계점}

본 연구는 다음과 같은 한계를 가진다.

첫째, 인프라 함정 카탈로그(F-1\textasciitilde F-6)는 N=1 프로젝트의 경험적 사례로, 다른 OQS 스택 구성이나 조직 환경에서 동일한 함정이 발생한다는 통계적 보장은 없다.

둘째, Stage~3는 순수 PQC-only TLS의 표준화 및 구현체 지원이 아직 성숙하지 않은 제약으로 인해, Advanced Hybrid PQC (NIST L5) 구성으로 검증하였다.
순수 PQC-only TLS는 관련 표준과 구현체 지원이 안정화된 이후의 향후 연구로 남긴다.

셋째, LLM 평가는 단일 모델(gpt-4o-mini)과 합성 패턴 N=16만을 대상으로 하며, GPT-4o·Claude·Gemini 등 더 큰 모델이나 실제 OSS 코드에서의 결과는 다를 수 있다.

넷째, 본 CBOM 모듈의 5 layer(자산 식별·3중 일관성 검증·진척도 판정·TLS 교차 검증·스키마 적합성)에 대한 누적 통계 기반 정량 평가는 본 연구 범위에서 수행하지 않았다.
본 연구는 모듈의 설계 및 구현을 제시하였으며, 실 운영 환경에서의 통과율 분포·MISMATCH 발생 빈도·회귀 감지 정확도 등의 측정은 별도 평가 환경 구축이 필요한 후속 과제다.

다섯째, 성능 측정은 GitHub Actions 러너와 로컬 Docker 환경에서 수행하였으므로 절대 성능값의 일반화에는 한계가 있다.
동일 환경 내 상대 비교로 해석해야 하며, 실 운영 트래픽 규모에서의 장기 안정성은 검증되지 않았다.

여섯째, 공인 CA 기반 PQC 인증서 체계가 적용되지 않았다.
Stage~3에서 사용한 ML-DSA-65 인증서는 실 공인 CA 체계가 아닌 프로젝트 내부 실험용 인증서로 구성하였으며, 실제 PKI 환경에서의 적용 가능성은 별도 검토가 필요하다.

\subsection{향후 연구 방향}

본 연구의 한계를 보완하기 위한 향후 연구 방향은 다음과 같다.

첫째, 순수 PQC-only TLS 지원이 성숙하는 시점에 Stage~3를 standalone ML-KEM 기반으로 완성하는 것이다.
IETF 표준화 완료 및 oqs-provider의 보안 강도 보고 정상화 이후 현재의 Advanced Hybrid 구성에서 PQC-only로의 전환이 가능할 것이다.

둘째, ML-DSA 기반 인증서 서명의 실 PKI 통합이다.
현재 프로젝트 내부 실험용 인증서 방식에서 내부 CA 기반 ML-DSA 인증서 발급 체계로 확장하면 실 운영 환경 적용 가능성을 높일 수 있다.

셋째, 고전 암호 탐지 시 자동 cipher replacement 기능의 강화다.
현재 파이프라인은 nginx TLS 설정의 자동 교체에 집중하지만, 코드베이스 내 레거시 암호 API를 자동 전환하는 LLM 마이그레이션 모듈의 정확도를 높이기 위해 더 큰 모델 적용 및 런타임 semantic test 통합이 필요하다.

넷째, 인프라 함정 카탈로그의 확장이다.
F-1\textasciitilde F-6는 단일 프로젝트 경험 기반이므로, BoringSSL, Java Bouncy Castle 등 다른 PQC 스택에서도 유사 패턴이 발생하는지 비교 연구를 통해 카탈로그를 확장할 수 있다.

다섯째, CBOM 모듈의 누적 통계 기반 정량 평가다.
본 연구에서 제안한 5 layer 검증 구조의 실효성을 다수 파이프라인 실행 데이터로 측정하고, 다양한 Stage 전환 시나리오에서 BASELINE·IMPROVED·UNCHANGED·REGRESSED 4상태 판정의 정확도, TLS 교차 검증의 PASS·MISMATCH·SKIPPED 분포, 3중 일관성 검증의 false positive·negative 비율 등을 통계적으로 분석할 필요가 있다.
또한 CBOM 범위를 현재의 TLS termination에서 데이터베이스 암호화·저장 데이터·내부 서비스 간 통신까지 확장하는 방향도 함께 검토할 수 있다.